\documentclass[../../../include/open-logic-section]{subfiles}

\begin{document}

\olfileid{sth}{z}{power}
\olsection{ઘાતગણો}

બીજી સ્વયંસિદ્ધિ સાથે આગળ વધીએ:

\begin{axiom}[ઘાતગણો]
કોઈ પણ ગણ $A$ માટે ગણ $\Pow{A} = \Setabs{x}{x \subseteq A}$
અસ્તિત્વ ધરાવે છે.
\[
\forall A \exists P \forall x(x \in P \liff (\forall z \in x)z \in A)
\]
\end{axiom}

આનું ઔચિત્ય બહુ સીધું છે. ધારો કે $A$ તબક્કા $S$ પર રચાય છે.
તો $A$ના બધા સભ્યો $S$ પહેલાં ઉપલબ્ધ હતા (\stagesacc પ્રમાણે).
તેથી પૃથક્કરણના ઔચિત્યમાં જે રીતે વિચાર્યું હતું એ જ રીતે
$A$નો દરેક ઉપગણ તબક્કા $S$ સુધીમાં રચાય છે. આમ તેઓ બધા
$S$ પછીના કોઈ પણ તબક્કે એક ગણ તરીકે રચાવા માટે ઉપલબ્ધ છે.
અને એવો કોઈ તબક્કો છે એ જાણીએ છીએ, કારણ કે $S$ અંતિમ તબક્કો
નથી (\stagessucc પ્રમાણે). તેથી $\Pow{A}$ અસ્તિત્વ ધરાવે છે.

ઘાતગણોની સ્વયંસિદ્ધિનું એક સરસ પરિણામ આ છે:

\begin{prop}\label{thm:Products}
કોઈ પણ ગણો $A, B$ આપેલા હોય, તો તેમનો કાર્તેઝીય ગુણાકાર
$A \times B$ અસ્તિત્વ ધરાવે છે.
\end{prop}

\begin{proof}
ઘાતગણોની સ્વયંસિદ્ધિ અને
\olref[sth][z][pairs]{prop:pairsconsequences}થી ગણ
$\Pow{\Pow{A \cup B}}$ અસ્તિત્વ ધરાવે છે. તેથી પૃથક્કરણથી
આ ગણ અસ્તિત્વ ધરાવે છે:
\[
C = \Setabs{z \in \Pow{\Pow{A \cup B}}}{(\exists x \in A)(\exists y \in B) z = \tuple{x, y}}.
\]
હવે કોઈ પણ $x \in A$ અને $y \in B$ માટે
\olref[sth][z][pairs]{prop:pairsconsequences}થી ગણ
$\tuple{x, y}$ અસ્તિત્વ ધરાવે છે. વધુમાં $x, y
\in A \cup B$ હોવાથી $\{x\}, \{x, y\} \in \Pow{A \cup B}$
અને $\tuple{x,y} \in \Pow{\Pow{A \cup B}}$ મળે છે.
તેથી $A \times B = C$.
\end{proof}

આ સાબિતીમાં ઘાતગણ અને પૃથક્કરણ સાથે કામ કરે છે. અને તેમાં
આશ્ચર્ય નથી. પૃથક્કરણ વિના ઘાતગણોનો સિદ્ધાંત બહુ
\emph{શક્તિશાળી} ન હોત. છેવટે પૃથક્કરણ કહે છે કે ગણના કયા
ઉપગણો અસ્તિત્વ ધરાવે છે, અને તેથી દરેક ઘાતગણ કેટલો ``ભરેલો''
છે તે નક્કી કરે છે.

\begin{prob}
કોઈ પણ ગણો $A, B$ માટે બતાવો કે (i) પ્રદેશ $A$ અને વિસ્તાર
$B$ ધરાવતા બધા સંબંધોનો ગણ અસ્તિત્વ ધરાવે છે; અને
(ii) $A$થી $B$માં જતા બધા વિધેયોનો ગણ અસ્તિત્વ ધરાવે છે.
\end{prob}

\begin{prob}
$A$ને ગણ લો અને $\sim$ને $A$ પરનો સામ્ય સંબંધ લો.
$A$ પર $\sim$ મુજબના સામ્ય વર્ગોનો ગણ, એટલે કે
$\equivclass{A}{\sim}$, અસ્તિત્વ ધરાવે છે એ સાબિત કરો.
\end{prob}

\end{document}
